\documentclass[8pt]{beamer}
\beamertemplatetransparentcovereddynamic	
\usepackage{ucs}
\usepackage{color,soul}
\usetheme{Singapore}
\usecolortheme{default}
\usepackage{amsmath,amsfonts,amssymb,booktabs}
\usepackage{graphics}
\usepackage{hyperref}
\usepackage{sansmathaccent}
\usepackage[position=top]{subfig}
\usepackage{arydshln}
\usepackage{graphicx}
\usepackage{dsfont}
\usepackage{tikz}
\usepackage{comment}
\usepackage{pgfplots}
\usepackage{algorithm2e,algorithmic,adjustbox,bbm}
\pgfplotsset{compat=newest} 
\usetikzlibrary{pgfplots.groupplots}
\usetikzlibrary{positioning}

\usetikzlibrary{arrows,positioning,calc}
\tikzset{cause/.style = {->,> = stealth', very thick, shorten >= 2, shorten <= 2},
  error/.style = {->,> = stealth', , shorten >= 2}, var/.style = {rectangle,rounded
  corners,thick,draw=black!75,fill=black!5,align=center,inner sep=5} }

\pdfmapfile{+sansmathaccent.map}
\setbeamertemplate{footline}[frame number]
\beamertemplatenavigationsymbolsempty

%\setbeamercolor{block body}{bg=normal}
\setbeamercolor{block title}{bg=blue!60,fg=white}

%-------------------Titlepage content--------------------------------

\title{What drives demand for playlists on Spotify?}
\vspace{5 mm}
\author[shortname]{Max J. Pachali\inst{1}, Hannes Datta\inst{1}}
\institute[shortinst]{\inst{1} Tilburg University}

\date{
 Digital Marketing and Media Seminar\\
\today
%October 18, 2021 
}
%-------------------Document-------------------------------

\begin{document}
\frame{\titlepage}

%-------------------TOC----------------------------------------
%\section[Outline]{}
%\frame{\tableofcontents}



%%%%%%%%%%%%%%%%%%%%%%%%%%%%%%%%%%%%%%%%%%%%%%%%%%%%%%%%%%%%%%%
%-------------------Introduction------------------------------
%%%%%%%%%%%%%%%%%%%%%%%%%%%%%%%%%%%%%%%%%%%%%%%%%%%%%%%%%%%%%%%

\section{Introduction}

\frame{
\frametitle{Concerns about market power of digital platforms}

\begin{itemize}

	\item Today, online platforms play an important role in society and economy 

	\item Platforms have started to dominate many industries 
		\begin{itemize}
		\item online retailing (e.g., Amazon)
		\item search and search advertising (e.g., Google)
		\item social media (e.g., Facebook)

	\end{itemize}
	
	\item Policy-makers around the world call for heightened transparency %and antitrust enforcement
		\begin{itemize}
		\item U.S. Justice Department accuses Google of "illegally protecting its monopoly over search and search advertising" (20 Oct. 2020)\footnote{"U.S. Accuses Google of Illegally Portecting Monopoly", 20 Oct. 2020, https://www.nytimes.com/2020/10/20/technology/google-antitrust.html}
		\item U.S. Congressional Hearing: Amazon, Apple, Facebook and Google wield market power "to crush competitors and amass data, customers and sky-high profits" (30 Jul 2020)\footnote{"Amazon, Apple, Facebook and Google grilled on Capitol Hill over their market power", 30 Jul 2020, https://www.washingtonpost.com/technology/2020/07/29/apple-google-facebook-amazon-congress-hearing/}
			
		\item EU Observatory on the Online Platform Economy\footnote{https://ec.europa.eu/digital-single-market/en/policies/online-platforms} 
		\end{itemize}
	
	
\end{itemize}
	

}

%%%%%%%%%%%%%%%%%%%%%%%%%%%%%%%%%%%%%%%%%%%%%%%%%%%%%%%%%%%%%%%
%%%%%%%%%%%%%%%%%%%%%%%%%%%%%%%%%%%%%%%%%%%%%%%%%%%%%%%%%%%%%%%
%%%%%%%%%%%%%%%%%%%%%%%%%%%%%%%%%%%%%%%%%%%%%%%%%%%%%%%%%%%%%%%


\frame{
\frametitle{Platform power in the music streaming market}

\begin{itemize}

	\item Platforms have overtaken traditional revenue streams of content producers\footnote{In 2019, 80\% of content producer's revenues generated through online streaming.}
	\item In the digital age, users mostly consume music through listening to platform-curated playlists
	\item Concerns about the market power of music streaming services started after the rise of playlists on Spotify --- the leading service\footnote{Market shares: Spotify: $36$\%, Apple Music: $18$\%, Amazon Music $13$\%, Tencent Music: $10$\%, Google/YouTube Music $5$\%, see https://musically.com/2019/12/09/report-spotify-has-36-market-share-of-music-streaming-subs/.} 
	\item "Playlists are the top tool Spotify is currently employing to expand their platform empire"\footnote{Quote by Liz Pelly (\url{https://tisch.nyu.edu/about/directory/clive-davis-institute/elizabeth-pelly}) taken from \url{https://watt.cashmusic.org/writing/thesecretlivesofplaylists}, accessed October, 19, 2020.}
	\item Yet, not much is known about the drivers of users' playlist consumption on Spotify 
		
%		\item Our paper empirically investigates the relative power balance in the music streaming market using daily data on 1.2m playlists on Spotify (the dominant platform in this market)\footnote{Market shares: Spotify: $36$\%, Apple Music: $18$\%, Amazon Music $13$\%, Tencent Music: $10$\%, Google/YouTube Music $5$\%, see https://musically.com/2019/12/09/report-spotify-has-36-market-share-of-music-streaming-subs/.}
	
\end{itemize}
}

%%%%%%%%%%%%%%%%%%%%%%%%%%%%%%%%%%%%%%%%%%%%%%%%%%%%%%%%%%%%%%%
%%%%%%%%%%%%%%%%%%%%%%%%%%%%%%%%%%%%%%%%%%%%%%%%%%%%%%%%%%%%%%%
%%%%%%%%%%%%%%%%%%%%%%%%%%%%%%%%%%%%%%%%%%%%%%%%%%%%%%%%%%%%%%%


\frame{
\frametitle{Power (im-)balances in the music streaming market}

\begin{itemize}
	\item Understanding the drivers of playlist demand is important  
	\item If users mainly listen to a playlist because \dots	
	\begin{itemize}
 	\item[$\rightarrow$]  it contains songs of popular artists, not Spotify would be powerful, but major record labels that have signed popular artists
	\item[$\rightarrow$]  of the theme, mood, music topic, or the overall acoustic styles it is representing ("playlist brand equity"), curators would be powerful through the decisions of which artists to include and "support"
	\end{itemize} 
	
	\item Spotify also features a selected set of playlists on its app%, including the web player
	\begin{itemize}
 	\item[$\rightarrow$] If consumers would be very responsive to these "promotions" of playlists, it would equip Spotify with additional power
	\end{itemize} 

\end{itemize}
}


%%%%%%%%%%%%%%%%%%%%%%%%%%%%%%%%%%%%%%%%%%%%%%%%%%%%%%%%%%%%%%%
%%%%%%%%%%%%%%%%%%%%%%%%%%%%%%%%%%%%%%%%%%%%%%%%%%%%%%%%%%%%%%%
%%%%%%%%%%%%%%%%%%%%%%%%%%%%%%%%%%%%%%%%%%%%%%%%%%%%%%%%%%%%%%%


\frame{
\frametitle{Goal of our paper and data}

\begin{itemize}

\item We empirically quantify the main drivers of playlist choice on Spotify in a structural demand model
\item Our analysis relies on a unique daily data set of initially 1.2 million playlists in the period of October 2019 until March 2020 collected from \textit{Chartmetric.com}
\item The playlist data includes (amongst others):
\begin{itemize}
\item[$\rightarrow$] daily aggregate playlist \textbf{listeners}   
\item[$\rightarrow$] share of content owned by the three \textbf{major labels} (Universal, Sony and Warner) on playlists, 
\item[$\rightarrow$] proxy for the \textbf{content popularity} of songs on playlists, 
\item[$\rightarrow$] total duration of the playlist
\item[$\rightarrow$] \dots
\item[$\rightarrow$] \dots
\end{itemize}
\item We complement the playlist data with daily information about which playlists get featured on Spotify's app (scraped from \textit{everynoise.com/worldbrowser.cgi}) 	

\end{itemize}
}

%%%%%%%%%%%%%%%%%%%%%%%%%%%%%%%%%%%%%%%%%%%%%%%%%%%%%%%%%%%%%%%
%%%%%%%%%%%%%%%%%%%%%%%%%%%%%%%%%%%%%%%%%%%%%%%%%%%%%%%%%%%%%%%
%%%%%%%%%%%%%%%%%%%%%%%%%%%%%%%%%%%%%%%%%%%%%%%%%%%%%%%%%%%%%%%


\frame{
\frametitle{Main factors of playlist demand}

\begin{enumerate}
\item Persistent (time-invariant) preferences for Spotify playlists %Users mainly value Spotify playlists because of their attractive brand images created by professional editors who design a hypothesis and theme for every playlist
\begin{itemize}
\item[$\rightarrow$] Market share of platform-curated playlists would drop by 50\% if replacing users' preferences for Spotify playlists by users' average tastes for non-Spotify playlists
\item[$\rightarrow$] Strong brands of Spotify playlists formed by professional editors who design a hypothesis and theme for every playlist (+ data advantage)
\end{itemize}
\item Responsiveness to featured playlists in the "Genres \& Moods" browse sections in Spotify's app
\begin{itemize}
\item[$\rightarrow$]  Market share of Spotify playlists would drop by 18\% if the platform wouldn't highlight a very selective set of its playlists to users in the app
\end{itemize}
\item Valuation of the popularity of added music contents 
\begin{itemize}
\item[$\rightarrow$]
Due to the strong persistence of user preferences, Spotify playlists would only lose 10\% (9\%, 8\%) of market share if Universal (Sony, Warner) content would be removed
\end{itemize}
\end{enumerate}



}

%%%%%%%%%%%%%%%%%%%%%%%%%%%%%%%%%%%%%%%%%%%%%%%%%%%%%%%%%%%%%%%
%%%%%%%%%%%%%%%%%%%%%%%%%%%%%%%%%%%%%%%%%%%%%%%%%%%%%%%%%%%%%%%
%%%%%%%%%%%%%%%%%%%%%%%%%%%%%%%%%%%%%%%%%%%%%%%%%%%%%%%%%%%%%%%


%\frame{
%\frametitle{Simulated magnitudes of main demand factors}
%
%\begin{enumerate}
%\item Market share of platform-curated playlists would drop by 50\% after replacing estimated preferences for Spotify playlists by users' average tastes of non-Spotify playlists 
%\item Spotify playlists would lose about 18\% of their market share if the platform wouldn't highlight a selective set of their playlists to users in the app 
%\item Due to the strong persistence of user preferences, Spotify playlists would only lose 10\% (9\%, 8\%) of market share if Universal (Sony, Warner) content would be removed from their playlists
%\end{enumerate}

%}

%%%%%%%%%%%%%%%%%%%%%%%%%%%%%%%%%%%%%%%%%%%%%%%%%%%%%%%%%%%%%%%
%%%%%%%%%%%%%%%%%%%%%%%%%%%%%%%%%%%%%%%%%%%%%%%%%%%%%%%%%%%%%%%
%%%%%%%%%%%%%%%%%%%%%%%%%%%%%%%%%%%%%%%%%%%%%%%%%%%%%%%%%%%%%%%


\frame{
\frametitle{Related literature}

\begin{enumerate}

\item Spotify's playlist power: Focusing on about 10 exceptionally popular playlists, Aguiar and Waldfogel (2018) document Spotify's power  to  influence  song  and  artist  success%.   According  to  their  estimates,  appearing  on "Today’s Top Hits"  raises a song’s streams by almost 20 million
\begin{itemize}
\item[$\rightarrow$]  We decompose the main drivers of user demand for playlists using a representative sample of thousands of popular playlists
\item[$\rightarrow$] Our data contains daily information about which playlists get featured in Spotify's "Genres \& Moods" browse sections and we find that users are quite responsive to these "promotions" of playlists
\item[$\rightarrow$] Our findings shed light on the power imbalances between major music producers and Spotify 
\end{itemize}

\item Literature on media power:  Simonov and Rao (2021) estimate persistent preferences for state-owned media outlets in Russia that are the main reason for media power and censored reporting 
\begin{itemize}
\item[$\rightarrow$] In our setting, users' have strong tastes for platform-curated playlists that are the main driver of Spotify's power over music producers 
\end{itemize}

\end{enumerate}

}


%%%%%%%%%%%%%%%%%%%%%%%%%%%%%%%%%%%%%%%%%%%%%%%%%%%%%%%%%%%%%%%
%%%%%%%%%%%%%%%%%%%%%%%%%%%%%%%%%%%%%%%%%%%%%%%%%%%%%%%%%%%%%%%
%%%%%%%%%%%%%%%%%%%%%%%%%%%%%%%%%%%%%%%%%%%%%%%%%%%%%%%%%%%%%%%

\section{Data sources}

\frame{
\frametitle{Focus on the playlist ecosystem at Spotify}
%\framesubtitle{Observing Attention Power of Platform and Content Producers}
\vspace{0.1cm}
\begin{center}
\centering
\includegraphics[height=2.8cm]{input/playlists_wide.png}
\end{center}
\begin{itemize}
	%\item "\textit{Playlists are the top tool Spotify are currently employing to expand their platform empire}"%\footnote{Quote by Liz Pelly from \url{https://watt.cashmusic.org/writing/thesecretlivesofplaylists}.}
	\item Playlists are curated list of songs that are typically centered around a particular theme e.g., by genre ("Deep House"), mood ("Relax"), or activity 
	\item Mean (median) play duration of playlists equal 8h (4.6h)
	\item Millions of lists available, curated by Spotify, major and indie labels and users
	\item Most of Spotify's playlists are managed by professional editors that change compositions (i.e., add or remove tracks) on fixed intervals --- mostly at the weekly level
	 %\item The professionally managed playlists create their own brand --- each has its own hypothesis, theme or audience that editors were thinking about
	 \item Playlist editors select songs based on the perceived fit with the brand image of a playlist
	%\item Content producers generate revenues per stream of their tracks on a playlist\footnote{Consumption on playlists account for a $30$\% share of listening time, and therefore royalties. Another $30$\% comes from users' personal music libraries - likely \textit{resulting} from discovery via playlists.}  
	%\footnote{"[...] awesome way to reach huge audiences instantly [...]" (Avicii, 2018, https://tinyurl.com/yb7q2svn)} %https://open.spotify.com/track/7tJshy0IW7NZeBqrKSFJra
%	\item Spotify owns most of the highly popular playlists, followed by the three major music labels
%	\item Any preferential treatment of major labels compared to independent music labels (controlling for popularity/quality of content) is potentially signal for major-label power over Spotify
\end{itemize}
}

%%%%%%%%%%%%%%%%%%%%%%%%%%%%%%%%%%%%%%%%%%%%%%%%%%%%%%%%%%%%%%%
%%%%%%%%%%%%%%%%%%%%%%%%%%%%%%%%%%%%%%%%%%%%%%%%%%%%%%%%%%%%%%%
%%%%%%%%%%%%%%%%%%%%%%%%%%%%%%%%%%%%%%%%%%%%%%%%%%%%%%%%%%%%%%%

\frame{
\frametitle{The Spotify playlist landscape}

%\begin{itemize}
%	\item There millions of lists available on Spotify every day that can be summarized into the following broad groups:
%\end{itemize}

\begin{enumerate}
\item \textbf{Spotify editorial playlists}
	\begin{itemize}
	\item{Like "Deep House Relax"}
	\item{Curated by a team of music experts and genre specialists from around the world hired by Spotify}
	\item{Spotify curates about 14,961 playlists and attracts $63$\% of all playlist followers}	  
	\end{itemize} 
\item \textbf{Professionally curated playlists}
	\begin{itemize}
	\item{The three major music labels and many independent labels also operate own playlist brands}\footnote{Sony (with Filtr), Universal (with Digstr), and Warner Music (with Topsify).}
	\item{"Tastemakers" --- artists, radio stations, brands, or influencers that curate music}
	\end{itemize} 
	\item \textbf{Spotify algorithmic playlists}
	\begin{itemize}
	\item Automatically created for each user by Spotify's own software algorithms
	\item Users access these recommendations in the "Made for You" section in the Spotify app (e.g., "Discover Weekly", "Release Radar", "On Repeat", "Your Daily Mixes", "Family Mix")
	\item "Discover Weekly" and "Release Radar" are updated weekly on every Monday and every Friday respectively
	\end{itemize}
	 \item \textbf{User playlists}
	\begin{itemize}
	\item Users can also share their playlists with one another and, in doing so, generate "followers" for their own libraries
	\end{itemize}

\end{enumerate}
}

%%%%%%%%%%%%%%%%%%%%%%%%%%%%%%%%%%%%%%%%%%%%%%%%%%%%%%%%%%%%%%%%%%%%%%%%%%%%
%-------------------Data---------------------------------------------
%%%%%%%%%%%%%%%%%%%%%%%%%%%%%%%%%%%%%%%%%%%%%%%%%%%%%%%%%%%%%%%%%%%%%%%%%%%%

\frame{
\frametitle{Data}
\framesubtitle{Playlist data}

%\begin{center}
%\centering
%\includegraphics[height=1cm]{input/chartmetric_logo.png}
%\end{center}

\begin{itemize}
	
	\item 1.2m playlists and daily data obtained from \textit{Chartmetric.com} (October 2019-March 2020) 
	\item Follower distribution across playlists is very skewed: the top 1\% (12,153) of playlists add up to a follower share of 84\% of the 1.2 million lists
	\item About 50\% of all playlists tracked have 1 follower (i.e., they are not shared with other users)
	%\item We classify the various unique playlist owners into the following groups:
	
\end{itemize}
		
\begin{table}


\centering
\begin{tabular}[t]{lrrr}
\hline \hline
Playlist owner & Follower share & N & Mean followers\\
\hline
Spotify & 0.633 & 14,961 & 108,506\\
Professional curators & 0.122 & 69,460 & 4,937\\
other & 0.101 & 352,940 & 743\\
Artists & 0.064 & 86,090 & 1,965\\
Major label (Sony) & 0.027 & 4,591 & 14,643\\
\addlinespace
Independent labels & 0.022 & 31,575 & 1,824\\
Major label (Universal) & 0.014 & 4,676 & 7,441\\
Major label (Warner) & 0.010 & 3,172 & 7,502\\
Users & 0.007 & 647,835 & 31\\
\hline \hline
\multicolumn{4}{l}{\textsuperscript{1} Playlist owners shown in descending order of their follower share.}\\
\end{tabular}
%\caption*{Playlist ownership in full sample.}
\label{follower_share_curators}
\end{table}
	
}


%%%%%%%%%%%%%%%%%%%%%%%%%%%%%%%%%%%%%%%%%%%%%%%%%%%%%%%%%%%%%%%
%%%%%%%%%%%%%%%%%%%%%%%%%%%%%%%%%%%%%%%%%%%%%%%%%%%%%%%%%%%%%%%
%%%%%%%%%%%%%%%%%%%%%%%%%%%%%%%%%%%%%%%%%%%%%%%%%%%%%%%%%%%%%%%

\frame{
\frametitle{Data}
\framesubtitle{Playlist data over time}

	\begin{itemize}
	\item Daily variation of playlist characteristics
	\begin{itemize}
		\item[$\rightarrow$] share of content owned by the three major labels (Universal, Sony, Warner)%\footnote{We classify labels into their (major-label) parent brands using a regular-expression based classifier validated on a secondary dataset from industry. The algorithm is available as an R package at https://github.com/hannesdatta/musicMetadata.}
		\item[$\rightarrow$] metric for the popularity of content on playlists: the number of tracks of a playlist added to other playlists at a given day%\footnote{This measure intuitively implies that the content of a playlist is more popular if many of its tracks get added to other playlists on the platform at a given day.}  
		\item[$\rightarrow$] number of tracks on playlists
		\item[$\rightarrow$] average time since release of tracks ("track age")
		\item[$\rightarrow$] average length of tracks (in minutes), 
		\item[$\rightarrow$] nine numeric measures characterizing the acoustic attributes on playlists (danceability, energy, speechiness, acousticness, instrumentalness, liveness, valence, tempo, loudness)
	\end{itemize}

	\item Daily variation of recommendation intensity at the platform level
	\begin{itemize}
		\item[$\rightarrow$] proxied by the daily refreshment\footnote{Defined as the number of songs added to personalized playlists divided by the number of personalized playlists in our sample.} of personalized playlists in our sample
	\end{itemize}
	
	\item \textit{Chartmetric.com}'s estimate of aggregate daily playlist listeners\footnote{The company collects the number of listeners of artists added to the playlists. Daily listeners are then computed as the mean number of listeners of (observable) artists added to the playlist.}	
	\end{itemize}
	
}

%%%%%%%%%%%%%%%%%%%%%%%%%%%%%%%%%%%%%%%%%%%%%%%%%%%%%%%%%%%%%%%
%%%%%%%%%%%%%%%%%%%%%%%%%%%%%%%%%%%%%%%%%%%%%%%%%%%%%%%%%%%%%%%
%%%%%%%%%%%%%%%%%%%%%%%%%%%%%%%%%%%%%%%%%%%%%%%%%%%%%%%%%%%%%%%

 \frame{
 \frametitle{Data}
 \framesubtitle{Merge feature data to playlist data}
\begin{columns}
\begin{column}{0.6\textwidth}
  
 \begin{itemize}
 	\item Spotify features a selected number of playlists in their app (called "Genres \& Moods" on Spotify)
 	\begin{itemize}
 				\item[$\rightarrow$] by country and categories (e.g., "mood", "dinner", "pop", ''EDM'')
				%\item[$\rightarrow$] most categories pertain to established genres (e.g., ''pop'', ''rock'', ''jazz'', ''EDM''). Other categories are rather new, like "sleep" or "workout" music
				 %\item[$\rightarrow$] featured lists are not personalized within countries

 	\end{itemize}

 	\item We retrieve daily data on featured playlists on Spotify from \textit{Everynoise.com/worldbrowser.cgi} (October 2019 until March 2020)
 	\begin{itemize}
 		\item data records a time stamp, the name of the section (e.g., ''latin'', ''focus'', ''chill'', ''rock'', etc.), the country name, and a list of playlist IDs featured in the category/country at the time of data retrieval
 		\item 6,739 playlists were featured during our period of observations in any of the sections and hence a very small fraction (0.56\%) of all playlists
 	\end{itemize}
 	\item Merge data on featured playlists with playlist data by IDs and dates 
 \end{itemize}


\end{column}
\begin{column}{0.4\textwidth}  %%<--- here
\begin{center}
\centering
\includegraphics[height=4cm]{input/promotions_spotify.png}
\end{center}
\end{column}
\end{columns}

%\footnotetext[1]{https://github.com/hannesdatta/data-spotify-promotions-releases}

}

%%%%%%%%%%%%%%%%%%%%%%%%%%%%%%%%%%%%%%%%%%%%%%%%%%%%%%%%%%%%%%%
%%%%%%%%%%%%%%%%%%%%%%%%%%%%%%%%%%%%%%%%%%%%%%%%%%%%%%%%%%%%%%%
%%%%%%%%%%%%%%%%%%%%%%%%%%%%%%%%%%%%%%%%%%%%%%%%%%%%%%%%%%%%%%%

 \frame{
 \frametitle{Data}
 \framesubtitle{Categorization of playlist ecosystem}
  
 \begin{itemize}
 	
	\item We use up to eight tag words (keywords) associated with playlists in the data to classify playlists into music genres using an association rule mining machine learning algorithm (e.g., Hastie et al.,2009)
	\item Why do we need to this? Spotify only features less than 1\% of playlists in our sample in the "Genres \& Moods" browse sections $\rightarrow$ we don't observe a genre classification for the vast majority of playlists
	\item Our categorization builds on the notion that playlists with similar keywords belong to the same music genre and are hence considered as closer substitutes by users
	\item Among the 1.2 playlists, we drop 436,183 playlists that do not contain any tag word in our sample
	\item For the remaining 779,117 playlists, the average number of tag words is 2.66
	\item To train the algorithm, we use a sub-sample of 100,000 playlists and 1,712 keywords to classify the total sample of playlists into broader categories (e.g., "EDM", "pop", "rock \& pop", etc.)

 \end{itemize}
	 
}


%%%%%%%%%%%%%%%%%%%%%%%%%%%%%%%%%%%%%%%%%%%%%%%%%%%%%%%%%%%%%%%
%%%%%%%%%%%%%%%%%%%%%%%%%%%%%%%%%%%%%%%%%%%%%%%%%%%%%%%%%%%%%%%
%%%%%%%%%%%%%%%%%%%%%%%%%%%%%%%%%%%%%%%%%%%%%%%%%%%%%%%%%%%%%%%

 \frame{
 \frametitle{Association rule mining algorithm (e.g., Hastie et al.,2009)}
 \framesubtitle{Categorization of playlist ecosystem}
  
 \begin{itemize}

\item Span the $(N \times P)$ matrix denoted as $X$ with elements $x_{n,p} \in \left\{0,1\right\}$, indicating whether a playlist $n$ contains keyword $p$
\item The goal of the algorithm is to find joint values of $X$ that appear most frequently ("mode finding")
\item We seek to find subsets of the integers $\mathcal{K} \subset \left\{1,\dots,P\right\}$ such that 

\begin{align}
    Pr\left[\bigcap_{k \in \mathcal{K}} (X_k = 1)\right] = Pr\left[\prod_{k \in \mathcal{K}} X_k = 1\right]
\end{align}

is large, where $\mathcal{K}$ is called an item set representing playlist keywords in our application and $X_k$ refers to the $k$-th column of $X$

 \end{itemize}

}

%%%%%%%%%%%%%%%%%%%%%%%%%%%%%%%%%%%%%%%%%%%%%%%%%%%%%%%%%%%%%%%
%%%%%%%%%%%%%%%%%%%%%%%%%%%%%%%%%%%%%%%%%%%%%%%%%%%%%%%%%%%%%%%
%%%%%%%%%%%%%%%%%%%%%%%%%%%%%%%%%%%%%%%%%%%%%%%%%%%%%%%%%%%%%%%

 \frame{
 \frametitle{Association rule mining algorithm}
 \framesubtitle{Categorization of  playlist ecosystem}
  
 \begin{itemize}
 		
\item We seek all item sets $\mathcal{K}_l$ whose support exceeds the threshold $t = 0.1$: 

\begin{align}
\begin{split}
    \{\mathcal{K}_l | T(\mathcal{K}_l) & > t\},\ \text{where} \\ 
    T(\mathcal{K}_l) & = \hat{Pr}\left[\bigcap_{k \in \mathcal{K}_l} (X_k = 1)\right] = 
    \frac{1}{N} \sum_{i=1}^{N} \prod_{k \in \mathcal{K}_l} x_{ik}
\end{split}
\label{eq:rule_mining}
\end{align}

\item We thus focus on combination of keywords that co-occur in at least 10\% of playlists in our sample
\item This threshold seems small, but is not if considering the total number of keywords (1,712) and the average number of keywords per playlist (2.66) $\rightarrow$ 5,008,986,720 ways to pick three keywords
\item The final goal of the classification is to represent each higher support item set as a set of association rules: $$A \rightarrow B, \ \text{with}\ A \ \text{and}\ B\ \text{being disjoint subsets and}\ A  \cup B = \mathcal{K}$$ and $B$ is the "consequent" of the rule
\item Rule confidence: C = $P(A \bigcap B) / P(A)$. We set the confidence to 90\% 
\item Minimum lift threshold: C/T(B). We set the minimum threshold to 1

\end{itemize}
}


%%%%%%%%%%%%%%%%%%%%%%%%%%%%%%%%%%%%%%%%%%%%%%%%%%%%%%%%%%%%%%%
%%%%%%%%%%%%%%%%%%%%%%%%%%%%%%%%%%%%%%%%%%%%%%%%%%%%%%%%%%%%%%%
%%%%%%%%%%%%%%%%%%%%%%%%%%%%%%%%%%%%%%%%%%%%%%%%%%%%%%%%%%%%%%%

 \frame{
 \frametitle{Steps of algorithm}
 \framesubtitle{Categorization of  playlist ecosystem}
  
 \begin{itemize}

\item Start classification based on Spotify's pre-defined genres in the browse sections of featured playlists (31 main genres):
\vspace{0.1cm}

\textit{afro, alternative, anime,  arab, blues, chill, christian music,  classical, comedy, country, desi, electronic/dance, focus, folk, acoustic, funk, video game music, hip hop, indie, jazz, latin, metal, party, pop, punk, r\&b, regional mexican, rock, romance, sleep, soul, student}

\begin{enumerate}
\item For every $genre$ in $main\  genres$, select $\tilde{X} = \{X|genre\in keywords\}$, i.e. the subset of playlists that contain $genre$ as keyword
\item Apply the Apriori algorithm (Agrawal et al., 1995) to identify the higher support item sets (Equation \ref{eq:rule_mining}) from the subset of playlists that contain the main genre as keyword, i.e. $\mathcal{K}_l^g(\tilde{X})$ with $\tilde{X} = \{X| main\ genre\in keywords\}$ \footnote{By conditioning on a main genre at each iteration, we ensure that our classification is complete and consistent with Spotify's browse sections.} 
\item[$\rightarrow$] Apriori algorithm: \textbf{first pass} over the data computes the support of all single item sets; discard low support item sets; \textbf{the second pass} computes the support of all item sets of size two that can be formed from pairs of the single item sets surviving the first pass, on so forth 
\item Cast set of association rules surviving minimum confidence and lift:
			$$\mathcal{R}_g =\{A\rightarrow B | \mathcal{K}_l^g\}$$
\item Stack total rules  $\mathcal{R} =\{ \mathcal{R}_g\}$
\end{enumerate}

\end{itemize}

}

%%%%%%%%%%%%%%%%%%%%%%%%%%%%%%%%%%%%%%%%%%%%%%%%%%%%%%%%%%%%%%%
%%%%%%%%%%%%%%%%%%%%%%%%%%%%%%%%%%%%%%%%%%%%%%%%%%%%%%%%%%%%%%%
%%%%%%%%%%%%%%%%%%%%%%%%%%%%%%%%%%%%%%%%%%%%%%%%%%%%%%%%%%%%%%%

 \frame{
 \frametitle{Classification output}
 \framesubtitle{Categorization of  playlist ecosystem}
  
 \begin{itemize}
 		
\item We mine in total 145 rules for $A \rightarrow B, \  A  \cup B = \mathcal{K}$, requiring that the consequent (B) is a main genre 
\item Output structure: $keyword\ \rightarrow\ main\ genre$, e.g.: \\ $pop\ rap \ \rightarrow\ hip\ hop$\\  $neo\ soul \ \rightarrow \ r\&b$ \\ $vocal\ jazz\ \rightarrow\ romance$ 

\item A playlist may be associated with multiple main genres, i.e. some keywords may create an association with rock, while other keywords create an association with romance
\item Our final genre categorization takes associations to multiple main genres as separate clusters, i.e. rock \& romance is a different category than rock
\end{itemize}
}

%%%%%%%%%%%%%%%%%%%%%%%%%%%%%%%%%%%%%%%%%%%%%%%%%%%%%%%%%%%%%%%
%%%%%%%%%%%%%%%%%%%%%%%%%%%%%%%%%%%%%%%%%%%%%%%%%%%%%%%%%%%%%%%
%%%%%%%%%%%%%%%%%%%%%%%%%%%%%%%%%%%%%%%%%%%%%%%%%%%%%%%%%%%%%%%

 \frame{
 \frametitle{Top 20 categories derived from the association rule learning classification}
 \framesubtitle{Categorization of  playlist ecosystem}

\begin{table}
\begin{adjustbox}{max width=\textwidth}
\centering
\begin{tabular}{lrrl}
\hline \hline
Category name & N & Followers (in m.) & Top 3 playlists\\
\midrule
pop & 58,962 & 378.37 & Today's Top Hits, Global Top 50, Songs to Sing in the Car\\
latin & 17,394 & 174.60 & ¡Viva Latino!, Baila Reggaeton, Rock en Español\\
pop+rb+student & 73,123 & 130.09 & All Out 00s, Songs to Sing in the Shower, I Love My '00s R\&B\\
hiphop & 36,797 & 102.46 & RapCaviar, Get Turnt, I Love My '90s Hip-Hop\\
rock & 58,606 & 80.92 & Rock Classics, 90s Rock Anthems, Legendary\\
\addlinespace
chill+electronicdance+party & 19,755 & 73.62 & Beast Mode, Dance Party, Power Hour\\
folkacoustic+romance & 22,617 & 56.59 & Your Favorite Coffeehouse, Run Wild, Afternoon Acoustic\\
comedy+hiphop & 47,747 & 51.44 & This Is Drake, \#ThrowbackThursday, B.A.E.\\
chill+electronicdance+party+pop & 12,475 & 43.79 & mint, Motivation Mix, Dance Hits\\
rock+romance & 7,291 & 37.34 & All Out 80s, Have a Great Day!, All Out 70s\\
\addlinespace
electronicdance & 12,279 & 36.62 & mint Latin, Housewerk, Techno Bunker\\
country & 8,829 & 30.28 & Hot Country, Country Gold, Country Kind of Love\\
metal+party+rock & 14,366 & 30.05 & This Is Queen, Rock Party, ’80s Hard Rock\\
focus & 3,953 & 26.97 & Deep Focus, Intense Studying, Peaceful Guitar\\
comedy+hiphop+pop & 16,220 & 23.84 & Hot Rhythmic, Signed XOXO, Workout Twerkout\\
\addlinespace
regionalmexican & 2,439 & 22.71 & Los Que Mandan: Top 40, Corridos Perrones, Mal de Amores\\
classical & 21,130 & 18.52 & Classical Essentials, Classical New Releases, This Is Mozart\\
indie & 14,446 & 17.31 & Ultimate Indie, Feel-Good Indie Rock, All New Indie\\
electronicdance+pop & 6,130 & 17.27 & Pop Remix, Sommerhits, Partyhymnen\\
focus+sleep & 163 & 17.16 & Peaceful Piano, Sleep, Deep Sleep\\
\hline \hline
\end{tabular}
\end{adjustbox}
%\caption{Summary of the top 20 categories derived from the association rule mining of playlists' tag words.}
\label{top_categories}
\end{table}
}

%%%%%%%%%%%%%%%%%%%%%%%%%%%%%%%%%%%%%%%%%%%%%%%%%%%%%%%%%%%%%%%
%%%%%%%%%%%%%%%%%%%%%%%%%%%%%%%%%%%%%%%%%%%%%%%%%%%%%%%%%%%%%%%
%%%%%%%%%%%%%%%%%%%%%%%%%%%%%%%%%%%%%%%%%%%%%%%%%%%%%%%%%%%%%%%

 \frame{
 \frametitle{Data}
 \framesubtitle{Construction of estimation sample}
  
 \begin{itemize}

\item Our final estimation sample relies on 12,055 major playlists in 50 music categories covering 65\% of the total followers of the 779,117 playlists we were able to allocate into genre categories\footnote{We had to drop 436,183 playlists from our sample that did not contain any keyword.}
\begin{itemize}
 	\item[$\rightarrow$] About 50\% playlist were dropped that only had a single follower (no listener data available)
	\item[$\rightarrow$] Select the genre categories belonging the top 95\% of cumulative follower market share of categories with at least 50 playlists included per category
	\item[$\rightarrow$] Select the 95\% mostly followed playlists within considered categories 
	\item[$\rightarrow$] These steps reduce the number of considered playlists to 30,746 in 81 genre categories covering 88\% of followers of the 779,117 playlists
	\item[$\rightarrow$]  Among the 30,746 playlists, \textit{Chartmetric.com} only observes listener data for 12,055 playlists that define the estimation sample
 \end{itemize}
 \item We do not rely on the follower data as a metric for daily demand as users can passively follow many playlists without actually listening to them on a given day
 \end{itemize}

}

%%%%%%%%%%%%%%%%%%%%%%%%%%%%%%%%%%%%%%%%%%%%%%%%%%%%%%%%%%%%%%%
%%%%%%%%%%%%%%%%%%%%%%%%%%%%%%%%%%%%%%%%%%%%%%%%%%%%%%%%%%%%%%%
%%%%%%%%%%%%%%%%%%%%%%%%%%%%%%%%%%%%%%%%%%%%%%%%%%%%%%%%%%%%%%%


%\frame{
%\frametitle{Data}
%\framesubtitle{Observing Personalized Recommendations}
%\vspace{0.1cm}
%\begin{center}
%\centering
%\includegraphics[height=4cm]{input/made-for-you.png}
%\end{center}
%\vspace{0.0cm}
%\begin{itemize}
%	\item Spotify uses playlists to serve machine-generated, personalized recommendations
%	\item Users access these recommendations in the "Made for You" section in the Spotify App\footnote{Spotify also personalized content outside of the "Made for You" section, e.g., via semi-personalized/adaptive mass playlists. Unlike playlists in the "Made for You" section, users are unaware of these personalizations.}
%	\begin{itemize}
%	\item e.g., Discover Weekly, Release Radar, On Repeat, Repeat Rewind, Your Daily Mixes, Family Mix
%	\end{itemize}
	
	
%	\begin{itemize}
%		\item Uniquely Yours: On repeat ("The songs you can't get enough of right now"), Repeat Rewind ("Past songs that you couldn't get enough of"), Your Top Songs 20xx ("the songs you loved most this year"), Family Mix ("Listen together with the people on your family plan"), Your Summer Rewind ("featuring your old summer favorites")
%		\item Discover New Music: Release Radar ("Catch the latest music from artists you follow, plus new singles picked for you"), Discover Weekly ("Your weekly mixtape of fresh music. Enjoy new music and deep cuts picked for you")
%		\item Your Daily Mixes: up to six; new songs by known artists
%		\item Others: Tastebreakers\footnote{"Start 2019 by broadening your horizons - list of songs from genres and artists that you don't normally explore and we think you'll like"}, The Ones That Got Away\footnote{"[...] collection of songs you'll wish you'd have discovered earlier in the year"}
%		\end{itemize}
%\end{itemize}
%}


%%%%%%%%%%%%%%%%%%%%%%%%%%%%%%%%%%%%%%%%%%%%%%%%%%%%%%%%%%%%%%%%%%%%%%%%%%%%
%-------------------MODEL FREE EVIDENCE--------
%%%%%%%%%%%%%%%%%%%%%%%%%%%%%%%%%%%%%%%%%%%%%%%%%%%%%%%%%%%%%%%%%%%%%%%%%%%%

\section{Descriptives}

\frame{
\frametitle{Average share of major label content per curator type}


\begin{figure}
\centering
\includegraphics[scale=0.25]{../paper/input/figure_mlshares.pdf}
%\caption{Summary of the average share of major label content per curator type across the 12,055 playlists included in the estimation sample.}
\label{ml_by_curators}
\end{figure}
	
\begin{itemize}
	\item Labels push their own content (high ML share on major label lists, low ML share on independent label lists)
	\item Spotify lists have similar ML share as lists by professional curators, artists and "others" 
\end{itemize}

}

%%%%%%%%%%%%%%%%%%%%%%%%%%%%%%%%%%%%%%%%%%%%%%%%%%%%%%%%%%%%%%%
%%%%%%%%%%%%%%%%%%%%%%%%%%%%%%%%%%%%%%%%%%%%%%%%%%%%%%%%%%%%%%%
%%%%%%%%%%%%%%%%%%%%%%%%%%%%%%%%%%%%%%%%%%%%%%%%%%%%%%%%%%%%%%%

\frame{
\frametitle{Major labels own most popular content}
	
\begin{figure}[t!]
\centering
\includegraphics[scale=0.25]{../paper/input/figure_popularity_mlshare.pdf}
%\caption{Empirical relation between the share of major label content and our measure for playlist popularity ("the number of tracks of a playlist added to other playlists at any point in time").}
\label{ml_content_pop}
\end{figure}

\begin{itemize}
	\item Adding major label content to playlists is associated with higher overall content popularity
	%\item Satiation once ML share approaches $80$\%, robust across curator types
\end{itemize}
}


%%%%%%%%%%%%%%%%%%%%%%%%%%%%%%%%%%%%%%%%%%%%%%%%%%%%%%%%%%%%%%%
%%%%%%%%%%%%%%%%%%%%%%%%%%%%%%%%%%%%%%%%%%%%%%%%%%%%%%%%%%%%%%%
%%%%%%%%%%%%%%%%%%%%%%%%%%%%%%%%%%%%%%%%%%%%%%%%%%%%%%%%%%%%%%%

\frame{
\frametitle{Playlist listeners and major label share}
	
\begin{figure}[t!]
\centering
\includegraphics[scale=0.25]{../paper/input/figure_listeners_mlshare.pdf}
%\caption{Empirical relation between the share of major label content and playlist listeners.}
\label{dist_featured}
\end{figure}

\begin{itemize}
	\item Playlists with a larger major label share seem to have more listeners
\end{itemize}
}

%%%%%%%%%%%%%%%%%%%%%%%%%%%%%%%%%%%%%%%%%%%%%%%%%%%%%%%%%%%%%%%
%%%%%%%%%%%%%%%%%%%%%%%%%%%%%%%%%%%%%%%%%%%%%%%%%%%%%%%%%%%%%%%
%%%%%%%%%%%%%%%%%%%%%%%%%%%%%%%%%%%%%%%%%%%%%%%%%%%%%%%%%%%%%%%

\frame{
\frametitle{Over time variation on playlist level}
	
\begin{figure}[t!]
\centering
\includegraphics[scale=0.2]{../output/figures_and_tables/ml_share_pop_listeners_list1.pdf}
%\caption{Over time evolution of major label share, content popularity and daily listeners exemplified for the "Housewerk" playlist.}
\label{variation_playlists}
\end{figure}

\begin{itemize}
	\item The previous plots do not account for playlist-specific demand factors
	\item We rely on over time variation on the playlist level to identify demand effects
	\item Variations in major label share affect content popularity over time
	\item But also the outbreak of the Covid 19 pandemic reduced the popularity of house music tracks  
\end{itemize}
}


%%%%%%%%%%%%%%%%%%%%%%%%%%%%%%%%%%%%%%%%%%%%%%%%%%%%%%%%%%%%%%%
%%%%%%%%%%%%%%%%%%%%%%%%%%%%%%%%%%%%%%%%%%%%%%%%%%%%%%%%%%%%%%%
%%%%%%%%%%%%%%%%%%%%%%%%%%%%%%%%%%%%%%%%%%%%%%%%%%%%%%%%%%%%%%%

\frame{
\frametitle{Feature distribution across playlists}

\begin{figure}
\centering
\includegraphics[scale=0.25]{../paper/input/ndays_featured.pdf}
%\caption{Distribution of days on feature in at least one of the "Genres \& Moods" sections in at least one country across the 12,055 playlists in our estimation sample. Note: the vertical line depicts the maximum number of days in the sample.}
\label{days_featured_across_lists}
\end{figure}

\begin{itemize}
	\item The majority of playlists in our sample are never on feature (in at least one of "Genres \& Moods" sections in at least one country)
\end{itemize}

}

%%%%%%%%%%%%%%%%%%%%%%%%%%%%%%%%%%%%%%%%%%%%%%%%%%%%%%%%%%%%%%%
%%%%%%%%%%%%%%%%%%%%%%%%%%%%%%%%%%%%%%%%%%%%%%%%%%%%%%%%%%%%%%%
%%%%%%%%%%%%%%%%%%%%%%%%%%%%%%%%%%%%%%%%%%%%%%%%%%%%%%%%%%%%%%%

\frame{
\frametitle{Variation in featured playlists across curator types}

\begin{figure}[t!]
\centering
\includegraphics[scale=0.2]{../paper/input/feature_share_curators.pdf}
%\caption{Share of featured playlists in at least one of the "Genres \& Moods" sections in at least one country by different curator types. Note: the share of featured playlists curated by Spotify is implied as the residual share.}
\label{feature_share_curators}
\end{figure}

\begin{itemize}
	\item As of end November 2019 onward, Spotify decided to stop featuring major label playlist and substantially increased the support of its own playlists
\end{itemize}

}

%%%%%%%%%%%%%%%%%%%%%%%%%%%%%%%%%%%%%%%%%%%%%%%%%%%%%%%%%%%%%%%
%%%%%%%%%%%%%%%%%%%%%%%%%%%%%%%%%%%%%%%%%%%%%%%%%%%%%%%%%%%%%%%
%%%%%%%%%%%%%%%%%%%%%%%%%%%%%%%%%%%%%%%%%%%%%%%%%%%%%%%%%%%%%%%

\frame{
\frametitle{Decreasing major label share of featured songs}

\begin{figure}[t!]
\centering
\includegraphics[scale=0.15]{../paper/input/ml_share_featured_non_featured.pdf}
%\caption{Evolution of average major label share between featured and non-featured playlists over time. Note: the separating condition is being featured in at least one of the "Genres \& Moods" sections in at least one country at a given day.}
\label{ML_share_per_group}
\end{figure}

\begin{itemize}
	\item Regime shift induced a substantial reduction of featured tracks released by any of the three major labels
	\item Evidence for decreasing major label power on Spotify 
	\item Spotify aims to increase the market share of independent labels and pursues the "goal to incorporate new or smaller artists into a variety of playlists" (Spotify 2020)
\end{itemize}

}

%%%%%%%%%%%%%%%%%%%%%%%%%%%%%%%%%%%%%%%%%%%%%%%%%%%%%%%%%%%%%%%
%%%%%%%%%%%%%%%%%%%%%%%%%%%%%%%%%%%%%%%%%%%%%%%%%%%%%%%%%%%%%%%
%%%%%%%%%%%%%%%%%%%%%%%%%%%%%%%%%%%%%%%%%%%%%%%%%%%%%%%%%%%%%%%

\frame{
\frametitle{Listener gap between featured and non-featured playlists}

\begin{figure}[t!]
\centering
\includegraphics[scale=0.15]{../paper/input/listeners_featured_non_featured.pdf}
%\caption{Comparison of average daily listeners between featured and non-featured playlists over time. Note: the separating condition is being featured in at least one of the "Genres \& Moods" sections in at least one country at a given day.}
\label{listeners_per_group}
\end{figure}

\begin{itemize}
	\item Playlists being featured in at least one of the "Genres \& Moods" sections have a substantially higher number of daily listeners
\end{itemize}

}

%%%%%%%%%%%%%%%%%%%%%%%%%%%%%%%%%%%%%%%%%%%%%%%%%%%%%%%%%%%%%%%%%%%%%%%%%%%%
%-------------------Empirical Framework--------
%%%%%%%%%%%%%%%%%%%%%%%%%%%%%%%%%%%%%%%%%%%%%%%%%%%%%%%%%%%%%%%%%%%%%%%%%%%%

\section{Demand Framework}


%%%%%%%%%%%%%%%%%%%%%%%%%%%%%%%
%%%%%%%%%%%%%%%%%%%%%%%%%%%%%%%
%%%%%%%%%%%%%%%%%%%%%%%%%%%%%%%

\frame{
\frametitle{Modeling playlist "consumption"}

\begin{itemize}
%\item Moment of consumption is discrete-choice decision: Users can't listen to multiple playlists simultaneously
\item Playlists as "products" differentiated by acoustic attributes, content popularity, play duration, curator type, \dots
\item Every playlist has its own theme or hypothesis and users have different preferences for these "brands"     
\item Users can choose among thousands of differentiated playlists and we model this decision as a discrete-choice 
\end{itemize}

\begin{columns}
\begin{column}{0.23\textwidth}  %%<--- here
\begin{center}
\centering
\includegraphics[height=1.8cm]{input/listA}
\end{center}
\end{column}

\begin{column}{0.2\textwidth}  %%<--- here
\begin{center}
\centering
\includegraphics[height=1.8cm]{input/listB}
\end{center}
\end{column}

\begin{column}{0.3\textwidth}  %%<--- here
\begin{center}
\centering
\includegraphics[height=1.8cm]{input/listC}
\end{center}
\end{column}

\end{columns}
   

}
%%%%%%%%%%%%%%%%%%%%%%%%%%%%%%%
%%%%%%%%%%%%%%%%%%%%%%%%%%%%%%%
%%%%%%%%%%%%%%%%%%%%%%%%%%%%%%%

\frame{
\frametitle{Demand specification}

\begin{itemize}

\item At every day $t$, users choose between listening to one of the playlists on Spotify or consuming content at competing music streaming platforms
\item "The market": Total number of premium subscribers of seven leading music streaming platforms (Spotify, Deezer, Apple, Amazon, Tencent, Pandora, YouTube)
\item We neglect the impact of fees on platform choice (monthly fees are almost homogeneously set across platforms to \$10) \footnote{Monthly fees of premium subscriptions across platforms in 2020: Spotify \$9.99, Apply Music \$9.99, Deezer \$9.99, Amazon Music \$9.99 (None-Prime, without accounting for annual plan discounts), Pandora \$9.99, YouTube Music \$11.99.}  
\item Every day $t$, every consumer $i$ in the market chooses $p$ from the $P_t + 1$ options that maximizes utility\footnote{Consumers can switch platforms on a daily basis in our model and additional fees from temporal double subscriptions are "sunk costs".}  

\begin{align} 
\begin{split} 
u_{ipt}  =& x_{pt}^{T} \Delta + \xi_{pt} + \zeta_{igt} + (1- \sigma) \cdot \varepsilon_{ipt}, \\
\end{split}
\label{first_demand}
\end{align}

\item $x_{pt}^{T} \Delta$ is the mean utility component from playlist $p$ at $t$ (details on next slide)
\item $\xi_{pt}$ is an unobserved time-varying playlist popularity or quality term 
\item $\zeta_{igt}$ is a random utility term such that $\zeta_{igt} + (1- \sigma) \cdot \varepsilon_{ipt}$ has an extreme value distribution
\end{itemize}

}


%%%%%%%%%%%%%%%%%%%%%%%%%%%%%%%
%%%%%%%%%%%%%%%%%%%%%%%%%%%%%%%
%%%%%%%%%%%%%%%%%%%%%%%%%%%%%%%

\frame{
\frametitle{Empirical specification of consumer utility}

\begin{align} 
\begin{split} 
u_{ipt}  %=& \xi_{pt} + x_{pt}^{T} \Delta + \zeta_{igt} + (1- \sigma) \cdot \varepsilon_{ipt} \\
=&  \gamma_p + \eta^{feat} N_{pt} + \beta^{pop} \text{ln}\left(1+n^{\text{-}p}_{pt}\right) + \omega r_{t} - \alpha c_{pt} + \\ &   + \gamma_{M(t),C(p)} + \gamma_{D(t)} + \rho t + \kappa \mathbbm{1}\{t \geq Covid\} + \boldsymbol{\phi}^{T} \boldsymbol{x}_{pt} + \zeta_{igt} + (1- \sigma) \cdot \varepsilon_{ipt}
\end{split}
\end{align} 

\begin{itemize}
\item $\gamma_p$ denotes users' time-invariant preferences for playlist $p$
\item $\eta^{feat}$ measures the effect of featuring playlists in Spotify's "Genres \& Moods" browse sections; $N_{pt}$ is the number of sections playlist $p$ is featured in aggregated over countries 
\item $n^{\text{-}p}_{pt}$ proxies for "content popularity" using the number of tracks from the focal playlists $p$ added to other playlists $\text{-}p$ 
\item $\omega$ measures the effect of personalized recommendation intensity\footnote{Defined as the number of songs added to personalized playlists divided by the number of personalized playlists in our sample.}
\item $c_{pt}$ is the playlist duration (in hours) and $\alpha$ is the marginal utility of time
\item $\gamma_{M(t),C(p)}$ represent month-category dummies, $\gamma_{D(t)}$ is a day-of-the-week dummy, $\rho$ controls for a general trend ("platform growth") and $\kappa$ accounts for the outbreak of the Covid 19 pandemic
\item $\boldsymbol{\phi}^{T} \boldsymbol{x}_{pt}$ controls for track age as well as acoustic attributes of $p$ 
\item We partition playlists into $G=50$ nests (obtained from the machine learning playlist categorization) and $\sigma$ is the nesting parameter %measuring the degree of intragroup correlation in consumer preferences 

\end{itemize}

}


%%%%%%%%%%%%%%%%%%%%%%%%%%%%%%%
%%%%%%%%%%%%%%%%%%%%%%%%%%%%%%%
%%%%%%%%%%%%%%%%%%%%%%%%%%%%%%%

\frame{
\frametitle{Estimation}

\begin{itemize}
\item Omitting the $t$ index and setting $\delta_p := x_{pt}^{T} \Delta + \xi_{pt}$, the market share of playlist $p \in H_g$ is given as: 

\begin{align} 
s_p  = \frac{\text{exp}\left[\delta_p/(1-\sigma)\right]}{\left(\sum_{h \in H_g} \text{exp}\left[\delta_h/(1-\sigma)\right] \right)^{\sigma} \cdot \sum_g \left\{ \left(\sum_{h \in H_g} \text{exp}\left[\delta_h/(1-\sigma)\right] \right)^{1-\sigma} \right\}},
\end{align} 

\item where $H_g$ is the set of playlists within the nest $g$
\item We invert market shares as proposed in Berry (1994) and estimate the model via OLS:

\begin{align} 
\begin{split} 
\text{ln} s_{pt} - \text{ln} s_{0t}  =& \gamma_p + \eta^{feat} N_{pt} + \beta^{pop} \text{ln}\left(1+n^{\text{-}p}_{pt}\right) + \omega r_{t} + \sigma \text{ln} s_{pt|gt} - \alpha c_{pt} + \\ &   + \gamma_{M(t),C(p)} + \gamma_{D(t)} + \rho t + \kappa \mathbbm{1}\{t \geq Covid\} + \boldsymbol{\phi}^{T} \boldsymbol{x}_{pt} + \epsilon_{pt},
\label{esti_eq}
\end{split}
\end{align} 

\item $s_{pt|gt}$ represents market share of playlist $p$ in nest $g$ and $\epsilon_{pt}$ corresponds to residual factors affecting demand for playlist $p$ at day $t$

\end{itemize}

}

%%%%%%%%%%%%%%%%%%%%%%%%%%%%%%%
%%%%%%%%%%%%%%%%%%%%%%%%%%%%%%%
%%%%%%%%%%%%%%%%%%%%%%%%%%%%%%%

\frame{
\frametitle{Endogeneity}

\begin{itemize}

\item We use instruments to correct for the potential endogeneity in the timing of featuring playlists and within nest market share
\item \textbf{Playlist featuring}: Daily number of "Genres \& Moods" sections that very small playlists are featured in (belonging to the 10\% of observations with the lowest followers in the same category) 
\item Intuition: Simultaneous coordination of the decision which playlists to feature within categories increases cost efficiency (relevance)\footnote{First-stage F-statistic: 87.55.}, but very small lists unlikely affect demand of major lists (different markets: "mainstream" vs. "niche") 
\item \textbf{Within nest market share}: Sum of play duration (number of tracks $\times$ length in hours) of competing playlists at every day (BLP-type) 
\item Intuition: This instrument is exogeneous (i.e., unrelated to $\epsilon_{pt}$) because the expansion of tracks on the platform rather reflects supply and not demand shocks, but negatively correlates\footnote{First-stage F-statistic: 217.34.} with $s_{pt|gt}$
%\item We view the variation in remaining variables as quasi-exogenous because of the \textit{rich set of fixed effects} included in our demand model and the \textit{high frequency of our sampled data set} (e.g., Rossi 2018)\footnote{Industry insiders (Heuvelings, 2020) describe that playlist compositions get updated on fixed intervals (mostly on a weekly level). It is therefore unlikely that curators adjust playlist compositions based on daily unobservables.} %See: https://www.bol.com/nl/f/auxiety/9200000131139132/

\end{itemize}

}

%%%%%%%%%%%%%%%%%%%%%%%%%%%%%%%
%%%%%%%%%%%%%%%%%%%%%%%%%%%%%%%
%%%%%%%%%%%%%%%%%%%%%%%%%%%%%%%

\frame{
\frametitle{Endogeneity}

\begin{itemize}

\item Curators may change playlist compositions % (such as number of tracks, acoustic attributes, major label share and hence implied content popularity) 
 in response to daily shocks
\item However, industry insiders (e.g., Heuvelings, 2020) describe that playlist compositions get updated on fixed intervals on Spotify  --- mostly on a weekly level
\item The mean (median) average update frequency of playlists in our sample equals 11.8 (7.8) days

\begin{figure}[t!]
\centering
\includegraphics[scale=0.15]{../paper/input/update_frequency.pdf}
%\caption{Distribution of average update frequency (in days) across playlists in the estimation sample. Note: Update frequency is defined as the number of days until a new song has been added or an existing song was removed on playlists.}
\label{dist_update_freq}
\end{figure}

\item The majority of curators thus do not adjust playlist compositions based on daily unobservables and the lion's share of variation in our data holds strategic decisions of curators constant

\item We follow (Rossi 2018) who proposes to deal with potential endogenous variables by using data sampled at a much higher frequency than strategic decisions are made

%\item We view the variation in remaining variables as quasi-exogenous because of the \textit{rich set of fixed effects} included in our demand model and the \textit{high frequency of our sampled data set} (e.g., Rossi 2018)\footnote{Industry insiders (Heuvelings, 2020) describe that playlist compositions get updated on fixed intervals (mostly on a weekly level). It is therefore unlikely that curators adjust playlist compositions based on daily unobservables.} %See: https://www.bol.com/nl/f/auxiety/9200000131139132/

\end{itemize}

}


%%%%%%%%%%%%%%%%%%%%%%%%%%%%%%%%%%%%%%%%%%%%%%%%%%%%%%%%%%%%%%%%%%%%%%%%%%%%
%-------------------Results--------
%%%%%%%%%%%%%%%%%%%%%%%%%%%%%%%%%%%%%%%%%%%%%%%%%%%%%%%%%%%%%%%%%%%%%%%%%%%%

\section{Results}


\frame{
\frametitle{Nested logit regressions}

\begin{table}[t!] \centering 
\begingroup\fontsize{6pt}{5pt}\selectfont
\begin{tabular}{@{\extracolsep{1pt}}lcc} 
\\[-1.8ex]\hline 
\hline \\[-1.8ex] 
 & \multicolumn{2}{c}{Specifications} \\ 
\cline{2-3} 
\\[-1.8ex] & \multicolumn{2}{c}{} \\ 
 & OLS & IV \\ 
\\[-1.8ex] & (1) & (2)\\ 
\hline \\[-1.8ex] 
 \textbf{featured (\# browse sections)} & \textbf{0.001***} (0.0001) & \textbf{0.022***} (0.002) \\ 
  \textbf{content popularity} & \textbf{0.004***} (0.001) & \textbf{0.046***} (0.009) \\ 
  \textbf{recommendation intensity} & \textbf{0.001***} (0.0002) & \textbf{-0.0004*} (0.0002) \\ 
  \textbf{sigma} & \textbf{0.985***} (0.001) & \textbf{0.719***} (0.020) \\ 
  alpha & 0.0001*** (0.00003) & 0.002*** (0.0003) \\ 
  trend & -0.0001*** (0.00001) & -0.0003*** (0.00003) \\ 
  covid outbreak & -0.019*** (0.0004) & -0.007*** (0.002) \\ 
  Monday & -0.001*** (0.0002) & -0.003*** (0.0004) \\ 
  Tuesday & 0.004*** (0.0003) & -0.002*** (0.0004) \\ 
  Wednesday & 0.007*** (0.0003) & 0.001** (0.0004) \\ 
  Thursday & 0.003*** (0.0001) & -0.001* (0.0003) \\ 
  Friday & 0.003*** (0.0001) & 0.0005 (0.0004) \\ 
  Saturday & 0.004*** (0.0004) & -0.002*** (0.0004) \\ 
  ac - danceability & -0.030** (0.015) & -0.141 (0.148) \\ 
  ac - energy & 0.008 (0.018) & 0.166 (0.171) \\ 
  ac - speechiness & -0.013 (0.022) & -0.326 (0.203) \\ 
  ac - acousticness & 0.028*** (0.010) & 0.256** (0.115) \\ 
  ac - instrumentalness & -0.004 (0.013) & 0.013 (0.081) \\ 
  ac - liveness & -0.007 (0.013) & 0.191 (0.119) \\ 
  ac - valence & 0.024** (0.010) & 0.324*** (0.105) \\ 
  ac - tempo & 0.0001 (0.0001) & 0.0004 (0.001) \\ 
  ac - loudness & -0.0005 (0.001) & -0.008 (0.007) \\ 
 \hline \\[-1.8ex] 
Observations & 1,678,774 & 1,678,774 \\ 
R$^{2}$ & 0.982 & 0.965 \\ 
Adjusted R$^{2}$ & 0.982 & 0.965 \\ 
F Statistic & 33,282.300*** (df = 223; 1666496) & 25,386.850*** \\ 
\hline 
\hline \\[-1.8ex] 
\textit{Note:}  & \multicolumn{2}{r}{*$p<$0.1; **$p<$0.05; ***$p<$0.01} \\ 
 & \multicolumn{2}{r}{Robust SEs in parentheses. Table omits playlist and month-category FEs.} \\
\end{tabular}
\endgroup
%\caption{Playlist demand estimates.} 
\label{tab_demand_1}
\end{table}

%\begin{itemize}
%	\item Promotions of playlists increase user demand (IV corrects attenuation bias)
%	\item Content popularity of playlists positively affects user demand 
%	\item More recommendations increase user demand for major playlists on Spotify (complementarity effect) 
%\end{itemize}
	
}


%%%%%%%%%%%%%%%%%%%%%%%%%%%%%%%
%%%%%%%%%%%%%%%%%%%%%%%%%%%%%%%
%%%%%%%%%%%%%%%%%%%%%%%%%%%%%%%

\frame{
\frametitle{Estimates of playlist "brand" preferences}

\begin{figure}
\centering 
\includegraphics[scale=0.18]{../paper/input/fes_follower_by_type.pdf}
%\caption{Relation between estimates of time-invariant playlist preferences (y-axis) and the logarithm of playlist followers (in April 2020).}
\label{fes_per_group}
\end{figure}

}


%%%%%%%%%%%%%%%%%%%%%%%%%%%%%%%%%%%%%%%%%%%%%%%%%%%%%%%%%%%%%%%%%%%%%%%%%%%%
%-------------------Counterfactuals---------------------------------------------
%%%%%%%%%%%%%%%%%%%%%%%%%%%%%%%%%%%%%%%%%%%%%%%%%%%%%%%%%%%%%%%%%%%%%%%%%%%%

\section{Counterfactuals}


\frame{
\frametitle{Counterfactuals}

\begin{enumerate}
\item Compare the magnitudes of main factors of playlist demand: By how much would market share of Spotify-curated playlists drop \dots
	\vspace{0.1cm}
	\begin{itemize}
		\item[(a)]  \dots if they wouldn't be associated as such strong brands by users? 
			\vspace{0.15cm}
		\item[(b)] \dots if they wouldn't be featured in the "Genres \& Moods" sections?
			\vspace{0.15cm}
		\item[(c)]  \dots if they wouldn't contain any (popular) Universal, Sony or Warner tracks?
	\end{itemize}
	\vspace{0.25cm}
\ 
\item Illustrate dependency of major label revenues on Spotify's editorial decisions
	\vspace{0.1cm}
	\begin{itemize}
		\item[(a)] By how much would revenues of Universal, Sony and Warner drop would their content be removed from featured playlists in Spotify's browse sections? 
%			\vspace{0.15cm}
%		\item[(b)] 
	\end{itemize}
%	\vspace{0.25cm}
\end{enumerate}

}

%%%%%%%%%%%%%%%%%%%%%%%%%%%%%%%
%%%%%%%%%%%%%%%%%%%%%%%%%%%%%%%
%%%%%%%%%%%%%%%%%%%%%%%%%%%%%%%

\frame{
\frametitle{Impact of time-invariant playlist preferences}

\begin{itemize}
\item Changes in playlist market share if users have substantially lower time-invariant preferences for Spotify playlists 
\item Set equal to average user preferences for non-Spotify playlists
\end{itemize}

\begin{table}[t!]
\centering
\begingroup\fontsize{8pt}{9pt}\selectfont
\begin{tabular}{llcccc}
  \hline \hline
    & & \multicolumn{3}{c}{Lower preferences for Spotify playlists} \\  Curator & base & counter & $\Delta$ & 95\% CI \\ \hline
 Outside good & 68.06 & 72.93 & 4.87 & [3.30, 6.14]  \\ 
   Spotify & 28.43 & 14.04 & -14.39 & [-17.60, -10.09]  \\ 
   Professional curators & 1.02 & 3.96 & 2.94 & [2.15, 3.75]  \\ 
   Artists & 1.01 & 3.81 & 2.80 & [2.00, 3.64]  \\ 
   Major label (Sony) & 0.52 & 2.37 & 1.85 & [1.26, 2.44]  \\ 
   other & 0.43 & 1.18 & 0.75 & [0.52, 0.97]  \\ 
   Independent labels & 0.23 & 0.42 & 0.20 & [0.13, 0.25]  \\ 
   Major label (Universal) & 0.17 & 0.75 & 0.58 & [0.43, 0.73]  \\ 
   Major label (Warner) & 0.14 & 0.52 & 0.39 & [0.28, 0.49]  \\ 
   %Users & 0.00 & 0.00 & -0.00 & [-0.00, -0.00] & 0.00 & 0.00 & [0.00, 0.00] \\ 
   \hline \hline
\end{tabular}
\endgroup
%\caption{Changes in playlist market share per curator group in two scenarios: 
%(i) users have substantially lower time-invariant preferences for Spotify playlists (set equal to average user preferences for non-Spotify playlists) as well as (ii) no playlists get 
%                               featured in the browse sections.} 
\label{ms_sims_fe_prom}
\end{table}

		
}



%%%%%%%%%%%%%%%%%%%%%%%%%%%%%%%
%%%%%%%%%%%%%%%%%%%%%%%%%%%%%%%
%%%%%%%%%%%%%%%%%%%%%%%%%%%%%%%

\frame{
\frametitle{Impact of featuring playlists in Spotify's "Genres \& Moods" sections}

\begin{itemize}
\item Changes in playlist market share if Spotify playlists don't get 
                               featured in the browse sections 
\end{itemize}

\begin{table}[t!]
\centering
\begingroup\fontsize{8pt}{9pt}\selectfont
\begin{tabular}{llcccc}
  \hline \hline
    &  & \multicolumn{3}{c}{No Spotify playlists featured} \\  Curator & base & counter & $\Delta$ & 95\% CI\\ \hline
 Outside good & 68.06  & 71.26 & 3.20 & [1.97, 5.10] \\ 
   Spotify & 28.43  & 23.30 & -5.13 & [-7.65, -3.23] \\ 
   Professional curators & 1.02 & 1.57 & 0.55 & [0.34, 0.80] \\ 
   Artists & 1.01 & 1.53 & 0.52 & [0.32, 0.75] \\ 
   Major label (Sony) & 0.52  & 0.82 & 0.30 & [0.18, 0.44] \\ 
   other & 0.43  & 0.68 & 0.25 & [0.16, 0.36] \\ 
   Independent labels & 0.23  & 0.36 & 0.13 & [0.08, 0.20] \\ 
   Major label (Universal) & 0.17 & 0.29 & 0.12 & [0.08, 0.17] \\ 
   Major label (Warner) & 0.14 & 0.20 & 0.06 & [0.04, 0.10] \\ 
   %Users & 0.00 & 0.00 & -0.00 & [-0.00, -0.00] & 0.00 & 0.00 & [0.00, 0.00] \\ 
   \hline \hline
\end{tabular}
\endgroup
%\caption{Changes in playlist market share per curator group in two scenarios: 
%(i) users have substantially lower time-invariant preferences for Spotify playlists (set equal to average user preferences for non-Spotify playlists) as well as (ii) no playlists get 
%                               featured in the browse sections.} 
\label{ms_sims_fe_prom}
\end{table}

		
}


%%%%%%%%%%%%%%%%%%%%%%%%%%%%%%%
%%%%%%%%%%%%%%%%%%%%%%%%%%%%%%%
%%%%%%%%%%%%%%%%%%%%%%%%%%%%%%%


\frame{
\frametitle{Empirical relation between major label share and content popularity}

\begin{itemize}
\item  To what extend are Spotify playlists dependent on popular major label tracks? 
\item Predict counterfactual content popularity of a playlist absent major label tracks
\item This is a problem of prediction in the following general form:

\begin{align} 
\begin{split} 
\text{ln}\left(1+n^{\text{-}p}_{pt}\right) =& \tilde{\gamma}_p +  f(ms_{pt}^{l}) + %\eta_{q}^{promo} N_{pt} + \omega_{q} r_{t} + \\ &+   \tilde{\gamma}_{M(t)} + \tilde{\gamma}_{D(t)} + \rho_{q} t + \boldsymbol{\phi}_{q}^{T} \boldsymbol{x}_{pt}  + 
\tilde{\varepsilon}_{pt}, 
\end{split} 
\end{align} 

where $l \in \{ \text{Universal, Sony, Warner} \}$, $ms_{pt}^{l}$ is the major label share of label $l$ added to playlist $p$ at $t$ and we specify $f(ms_{pt}^{l}) = \beta_l \mathbbm{1}\{ ms_{pt}^{l} \neq 0 \} + \tilde{\beta}_{l} ms_{pt}^{l}$

\end{itemize}

\begin{table}[H] \centering 
\begingroup\fontsize{6pt}{7pt}\selectfont
\begin{tabular}{@{\extracolsep{1pt}}lc} 
\\[-1.8ex]\hline 
\hline \\[-1.8ex] 
 & \multicolumn{1}{c}{\textit{Dependent variable:}} \\ 
\cline{2-2} 
\\[-1.8ex] & Content popularity \\ 
\hline \\[-1.8ex] 
 content share: Universal & 2.294$^{***}$ (0.163) \\ 
  content share: Warner & 2.145$^{***}$ (0.159) \\ 
  content share: Sony & 2.121$^{***}$ (0.135) \\ 
  content share: Universal \textgreater  0 & 0.147$^{***}$ (0.024) \\ 
  content share: Warner \textgreater  0 & 0.190$^{***}$ (0.024) \\ 
  content share: Sony \textgreater  0 & 0.173$^{***}$ (0.021) \\ 
 \hline \\[-1.8ex] 
Observations & 1,678,774 \\
R$^{2}$ & 0.258 \\ 
Adjusted R$^{2}$ & 0.252 \\ 
F Statistic & 117.577$^{***}$ (df = 6; 1666713) \\ 
\hline 
\hline \\[-1.8ex] 
\textit{Note:}  & \multicolumn{1}{r}{$^{*}$p$<$0.1; $^{**}$p$<$0.05; $^{***}$p$<$0.01} \\ 
 & \multicolumn{1}{r}{Robust SEs in parentheses. Table omits playlist FEs.} \\ 
\end{tabular}
\endgroup
  %\caption{Empirical relation between share of major label content and playlists' content popularity.}
  \label{tab:pop_reg} 
\end{table}

}
%%%%%%%%%%%%%%%%%%%%%%%%%%%%%%%
%%%%%%%%%%%%%%%%%%%%%%%%%%%%%%%
%%%%%%%%%%%%%%%%%%%%%%%%%%%%%%%

\frame{
\frametitle{Impact of removing Universal content}

\begin{itemize}
\item Changes in playlist market share if Spotify playlists removed all tracks released by Universal  
\end{itemize}

\begin{table}[H]
\centering
\begingroup\fontsize{8pt}{9pt}\selectfont
\begin{tabular}{llcccc}
  \hline \hline
    & & \multicolumn{3}{c}{No Universal content on Spotify playlists}  \\  Curator & base & counter & $\Delta$ & 95\% CI \\ \hline
 Outside good & 68.06 & 70.74 & 2.68 & [1.58, 3.68]  \\ 
   Spotify & 28.43 & 25.72 & -2.71 & [-3.75, -1.59]  \\ 
   Professional curators & 1.02 & 0.99 & -0.03 & [-0.05, -0.01]  \\ 
   Artists & 1.01 & 1.07 & 0.06 & [0.03, 0.10]  \\ 
   Major label (Sony) & 0.52 & 0.53 & 0.01 & [-0.00, 0.03]  \\ 
   other & 0.43 & 0.42 & -0.00 & [-0.01, 0.00]  \\ 
   Independent labels & 0.23 & 0.21 & -0.01 & [-0.02, -0.01]  \\ 
   Major label (Universal) & 0.17 & 0.18 & 0.01 & [0.01, 0.03] \\ 
   Major label (Warner) & 0.14 & 0.14 & -0.00 & [-0.01, 0.00]  \\ 
    %Users & 0.00 & 0.00 & -0.00 & [-0.00, -0.00] & 0.00 & -0.00 & [-0.00, -0.00] \\ 
   \hline \hline
\end{tabular}
\endgroup
%\caption{Changes in playlist market share per curator group if Spotify playlists removed   
%(i) all tracks released by Universal as well as (ii) all tracks released by Sony.} 
\label{ms_sims_no_uni_sony}
\end{table}

\begin{itemize}
\item Corresponding changes of market share of Spotify playlists if Sony and Warner content would be removed: -2.47, -2.38  
\end{itemize}

		
}


%%%%%%%%%%%%%%%%%%%%%%%%%%%%%%%
%%%%%%%%%%%%%%%%%%%%%%%%%%%%%%%
%%%%%%%%%%%%%%%%%%%%%%%%%%%%%%%

\frame{
\frametitle{Make $f(ms_{pt}^{l})$ more flexible?}

\begin{itemize}

\item Functional relation between major label share and playlist content popularity likely more complex than our current specification
\item \dots 	

\end{itemize}


}





%%%%%%%%%%%%%%%%%%%%%%%%%%%%%%%
%%%%%%%%%%%%%%%%%%%%%%%%%%%%%%%
%%%%%%%%%%%%%%%%%%%%%%%%%%%%%%%

\frame{
\frametitle{Major label revenues and Spotify's editorial decisions}

	\begin{itemize}
	\item{Major label's $l \in \{ \text{Universal, Sony, Warner} \}$ revenues from editorial playlists on Spotify can be expressed as follows:

		\begin{align}
		R_{l} = \sum_{p \in P_{\text{-}0}} s_p \left(\delta_{p} \right) \cdot  MS^{l}_{p}  \cdot D  \cdot r_l,  
		\label{eq:profit_labels}
		\end{align}
		
		}
 	\item{$P_{\text{-}0}$ is the set of playlists on Spotify (excluding the outside good), $s_p$ is the market share of playlist $p$,  $MS_p^l$ are the number of songs released by major label $l$ on $p$, $D$ is the market size and $r_l$ is the per stream royalty payment of label $l$} 	
		
	\end{itemize}
		
}


%%%%%%%%%%%%%%%%%%%%%%%%%%%%%%%
%%%%%%%%%%%%%%%%%%%%%%%%%%%%%%%
%%%%%%%%%%%%%%%%%%%%%%%%%%%%%%%

\frame{
\frametitle{Major label revenues and Spotify's editorial decisions}

%\begin{itemize}
%\item{\%-change of producers' revenues on the platform in the three scenarios that Spotify playlists remove Universal, Sony and 	Warner tracks:}
%\end{itemize}
%
%\begin{table}[H]
%\centering
%\begingroup\fontsize{6pt}{7pt}\selectfont
%\begin{tabular}{lcccccc}
 % \hline
 %  & \multicolumn{2}{c}{No Universal content} & \multicolumn{2}{c}{No Sony content}
 %                          & \multicolumn{2}{c}{No Warner content} \\ & $\Delta\ in\ \%$ & 95\% CI & $\Delta\ in\ \%$ & 95\% CI & $\Delta\ in\ \%$ & 95\% CI\\ \hline
%Universal & -90.65 & [-93.46, -88.69] & -4.98 & [-7.72, -2.54] & -4.89 & [-7.64, -2.48] \\ 
%  Sony & -3.13 & [-5.84, -0.93] & -85.58 & [-89.48, -82.94] & -3.67 & [-6.19, -1.45] \\ 
%  Warner & -1.72 & [-4.59, 0.91] & -2.28 & [-5.00, 0.29] & -93.51 & [-95.54, -92.06] \\ 
%  Independent labels & 45.66 & [39.90, 52.14] & 29.93 & [24.34, 35.53] & 22.58 & [17.53, 27.56] \\ 
%   \hline
%\end{tabular}
%\endgroup
%\caption{Percentage change of major label revenues for the counterfactual scenarios of no major label content on Spotify playlists.} 
%\label{rev_labels_scens_II}
%\end{table}

\begin{itemize}
\item{\%-change of producers' revenues on Spotify in the three scenarios of no major label tracks on featured playlists}

\begin{table}[H]
\centering
\begingroup\fontsize{6pt}{7pt}\selectfont
\begin{tabular}{lcccccc}
  \hline \hline
   & \multicolumn{2}{c}{No Universal content} & \multicolumn{2}{c}{No Sony content} & \multicolumn{2}{c}{No Warner content} \\ & $\Delta\ in\ \%$ & 95\% CI & $\Delta\ in\ \%$ & 95\% CI & $\Delta\ in\ \%$ & 95\% CI\\ \hline
Universal & -70.14 & [-77.48, -64.62] & -5.83 & [-8.63, -3.23] & -5.61 & [-8.38, -3.07] \\ 
  Sony & -3.93 & [-6.60, -1.67] & -58.38 & [-68.86, -51.12] & -4.88 & [-7.41, -2.51] \\ 
  Warner & -2.56 & [-5.34, 0.03] & -3.50 & [-6.20, -0.97] & -61.84 & [-71.42, -54.75] \\ 
  %Independent labels & 32.67 & [25.97, 40.30] & 17.31 & [10.72, 24.49] & 11.90 & [6.22, 17.89] \\ 
   \hline \hline
\end{tabular}
\endgroup
%\caption{Percentage change of major and independent label revenues for the counterfactual scenarios of no major label
 %                     content on featured playlists in Spotify's browse sections.} 
\label{rev_labels_featured}
\end{table}


\end{itemize}

		
}



%%%%%%%%%%%%%%%%%%%%%%%%%%%%%%%%%%%%%%%%%%%%%%%%%%%%%%%%%%%%%%%%%%%%%%%%%%%%
%-------------------Conclusions---------------------------------------------
%%%%%%%%%%%%%%%%%%%%%%%%%%%%%%%%%%%%%%%%%%%%%%%%%%%%%%%%%%%%%%%%%%%%%%%%%%%%

\section{Conclusion}

\frame{
\frametitle{Conclusion}
\begin{itemize}
	\item  We study the main drivers of playlist demand on Spotify using a unique data set on daily listeners and characterizing attributes of about 12,000 playlists
	\item We find that users mainly chose Spotify playlists for time-invariant reasons, most likely because of their attractive brand images
	\item Spotify employs a large body of professional editors who help creating the own hypothesis and theme ("brand equity") of a playlist 
	\item Users are also quite responsive to featuring playlists in the platform's "Genres \& Moods" browse sections $\rightarrow$ another part of Spotify's power
	\item In relative terms, user demand is less responsive to fluctuations in the popularity of added tracks on playlists 
	
	%\item Implications:
	%\begin{itemize}
	%	\item[(1)] Revenues of major labels (and music producers more broadly) largely depend on Spotify's decisions about playlist compositions and featuring their tracks in the "Genres \& Moods" browse sections
	%	\item[(2)] To be become more independent, music producers should further invest in growing the brand equity of their own playlist channels  
	%\end{itemize}


\end{itemize}
}

%%%%%%%%%%%%%%%%%%%%%%%%%%%%%%%
%%%%%%%%%%%%%%%%%%%%%%%%%%%%%%%
%%%%%%%%%%%%%%%%%%%%%%%%%%%%%%%

\frame{
\frametitle{Implications}
\begin{itemize}
		
\item[(1)] Revenues of major labels (and music producers more broadly) largely depend on Spotify's decisions about playlist compositions and featuring playlists in the "Genres \& Moods" browse sections
\item[(2)] In relative terms, Spotify is much less dependent on the availability of popular major label tracks on its playlists
\item[(3)] To reduce Spotify's "playlist monopoly" and increase competition, music producers and other curators should further invest in growing both the quantity and brand equity of their own playlist channels on Spotify

\end{itemize}

}
	

%%%%%%%%%%%%%%%%%%%%%%%%%%%%%%%%%%%%%%%%%%%%%%%%%%%%%%%%%%%%%%%%%%%%%%%%%%%%
%-------------------Next steps---------------------------------------------
%%%%%%%%%%%%%%%%%%%%%%%%%%%%%%%%%%%%%%%%%%%%%%%%%%%%%%%%%%%%%%%%%%%%%%%%%%%%

%
%\section{Next steps}
%
%\frame{
%\frametitle{Next steps}
%\begin{itemize}
%	
%	\item{Robustness check: Conditional estimation sample on playlists curated by others (no profit interest, no strategic choice of label content)} 
%	\item{Potential IV-candidates: Number of followers of label artists in different social media channels; lagged label shares of playlists} 
%	\item Control for major label artists' social media scores
%%	\item Make listener-to-follower ratio dynamic
%%	\item Allow for category-level heterogeneity
%%	\item Independent labels
%	
%\end{itemize}
%}




%%%%%%%%%%%%%%%%%%%%%%%%%%%%%%%%%%%%%%%%%%%%%%%%%%%%%%%%%%%%%%%%%%%%%%%%%%%%
%-------------------Conclusions---------------------------------------------
%%%%%%%%%%%%%%%%%%%%%%%%%%%%%%%%%%%%%%%%%%%%%%%%%%%%%%%%%%%%%%%%%%%%%%%%%%%%

%\section{Conclusion}

%\frame{
%\frametitle{Conclusion}
%\begin{itemize}
%	\item 
%\end{itemize}
%}



\frame{
\vspace{5 mm}
\huge
\begin{center}
\textbf{Thank you for your attention!}
\end{center}
}


%%%%%%%%%%%%%%%%%%%%%%%%%%%%%%%%%%%%%%%%%%%%%%%%%%%%%%%%%%%%%%%%%%%%%%%%%%%%
%-------------------Appendix------------------------------------------------
%%%%%%%%%%%%%%%%%%%%%%%%%%%%%%%%%%%%%%%%%%%%%%%%%%%%%%%%%%%%%%%%%%%%%%%%%%%%

\section{Appendix}


\frame{
\frametitle{Business Model of Digital Music Platforms}
\framesubtitle{Spotify as a two-sided market}
%\begin{center}
%\centering
%\includegraphics[width=1.0\textwidth]{input/newsmedia.png}
%\end{center}
\begin{itemize}
	\item Spotify
	\begin{itemize}
		\item Market demand depends not only on consumer preferences, but also on search costs
		\item Platform lowers search costs and affects discovery (Aguiar and Waldfogel 2018), by means of playlists, recommendations, and curated artist profiles.
	\end{itemize}
	\item Record labels
	\begin{itemize}
	\vspace{0.1cm}
		\item Major record labels account for 70\% of global industry revenue (Universal 29.8\%, Warner 22.6\%, Sony 18.1\%), and provide most of the highly popular content to Spotify %; own ~$xx$\% of Spotify shares
		\item Independent record labels: highly fragmented (10,000+)
	\end{itemize}
	\item Revenue model\footnote{https://heroic.academy/artist-guide-spotify-playlist-royalties-verified-profiles/} 
	\begin{itemize}
		\item Consumers pay monthly subscription fee, and consume music via playlists ($30$\% of consumption) or personal music libraries ($70$\% of consumption)
		\item Spotify collects revenue from ad-based and premium tiers; deducts operational expenses and pays out about $52$\% to labels, \emph{relative to a label's number of streams} $\rightarrow$ labels compete for plays
	\end{itemize}
	
\end{itemize}
}


\frame{
\frametitle{Industry debates relationship between Spotify and the major labels}
\begin{itemize}
	\item Payout relative to plays on platform seems democratic and fair, but...
	\item Major labels own(ed) about $6$-$10$\% of Spotify shares, and major-label content is critical to the business model of Spotify ("can't do without Justin Bieber")\footnote{https://www.musicbusinessworldwide.com/one-reason-why-spotifys-deals-with-the-major-labels-rest-on-a-knife-edge/}
	\vspace{0.4cm}
	\item Consequently, major labels negotiated better deals with Spotify, compared to independent labels\footnote{Sony contract leak in 2015, see https://www.theverge.com/2015/5/19/8621581/sony-music-spotify-contract.}
	\begin{itemize}
		\item Higher-than-standard royalty rates (minimum per-stream payouts)
		\item Control over share of advertising space, playlist addition, and playlist promotion
	\end{itemize}
	
	\vspace{0.4cm}
	\item Public debate on (fairness of) Spotify's business model
	\begin{itemize}
		\item Shouldn't independent music artists earn the same, compared to major-label artists?
		\item Shouldn't Spotify try to become label itself to reduce major-label power?
		\item How powerful is Spotify in diverting attention (e.g., by introducing a new mood music category)?
	\end{itemize}	
\end{itemize}

}

%References
%\printbibliography
%\bibliography{../Literatur/bibliography}
%\bibliographystyle{chicago}



\end{document}



